\subsection{Simplifying a biclique}\label{sec:simplifyingbiclique}
In this section, we prove the following lemma. Recall that, from Section~\ref{sec:cleaning}, for positive integer $t,p$, $\mathcal{R}'(t)$ is the smallest positive integer such that any bipartite graph with a matching of size $\mathcal{R}'(t)$ either has an induced matching of size $t$ or a $t$-induced-biclique.
Furethermore $f(t,p) = 2^{\mathcal{R}'(2t)} f(t,p-1) + \mathcal{R}'(2t)$ 
when $p>1$ and $f(t,1)=t$. 

\digging*

The remainder of this section is dedicated to the proof of Lemma~\ref{lem:digging}. Recall that $A,B$ are ordered sets. Let $g(t) = f(t, \mathcal{R}(2t,2))$.
Further let $A=(a_1, \ldots, a_{g(t)})$ and $B=(b_1, \ldots, b_{g(t)})$.  

We begin by showing that, one can assume without loss of generality that $D$ is acyclic.
Suppose that $D$ is not acyclic. First observe that for any $a_i,a_j \in A$, such that $i \neq j$, $a_i$ and $a_j$ belong to different strongly connected components of $D$. Similarly, each vertex of $B$ belongs to a distinct strongly connected component of $D$. Further, for any $a_i,b_j$, $a_i$ and $b_j$ do not belong to the same strongly connected component of $D$. Indeed, as otherwise there is a $(b_j,a_i)$-path in $D$ contracting the definition of an $(A,B)$-biclique. Thus, each vertex in $\{a_1, \ldots, a_{g(t)}, b_1 , \ldots, b_{g(t)}\}$ belong to distinct strongly connected components of $D$.
Let $D'$ be obtained from $D$ by contracting each strongly connected component of $D$ into a single vertex. That is, $D'$ has a vertex for each strongly connected component of $D$ and for two vertices $u,v$ in $D'$, there is an edge from $u$ to $v$ in $D'$ if there is vertex in the strongly connected component of $D$ which was contracted onto $u$, that has an edge in $D$ to a vertex of the strongly connected component of $D$ which was contracted to $v$. Let us call a vertex of $D'$ $a_i$ if it is obtained by contracting a strongly connected component containing $a_i$. Similarly, let us call a vertex of $D'$ $b_i$ if it is obtained by contracting a strongly connected component containing $b_i$. It is easy to observe that $D'$ also contains an $(A,B)$-matching and $D'^{\star}$ contains an $(A,B)$-induced-biclique. Indeed, since an $(a_i,b_j)$-path in $D$ implies an $(a_i,b_j)$-path in $D'$ and no $(a_i,a_j)$-path (resp.~no $(b_i,b_j)$-path) in $D$ implies no $(a_i,a_j)$-path (resp.~no $(b_i,b_j)$-path) in $D'$. Further if $D'$ contains an $(a_i,b_i)$-path of length strictly greater than $1$ then so does $D$. Since $D'$ is obtained by contraction operation from $D$, we conclude that without loss of generality, we can assume for the rest of the section that the input graph $D$ in Lemma~\ref{lem:digging} is acyclic.

We say that an edge $e \in E(D) \setminus E(A,B)$ is {\em contraction-redundant} with respect to $(A,B)$ in $D$ if the following holds.

\begin{enumerate}
\item If $D/e$ is acyclic, 
\item all the edges in $\{(a_i,b_i) : i \in [g(t)]\}$ (that is the matching edges of $(A,B)$) are minimal edges in $D/e$, and
\item $(D/e)^{\star}$ has an $(A,B)$-induced-biclique..
\end{enumerate} 

Let $\Dcontr$ be the graph obtained from $D$ by repeatedly contracting the contraction-redundant edges with respect to $(A,B)$ until $\Dcontr$ has no such edge. 
For the ease of notation by $aa$-path we mean a path between two distinct vertices of $A$, by $bb$-path we mean a path between two distinct vertices of $B$ and by $a_{\ast}b_{\ast}$-path we mean a $(a_i,b_i)$-path for some $i \in [g(t)]$.
We divide the proof of Lemma~\ref{lem:digging} into two independent parts based on the length of a longest path in $\Dcontr$. In Section~\ref{sec:biclique:longpath}, we consider the case when the length of a longest path in $D_{contr}$ is at least $2t+2$.
In this case, we show that $\Dcontr$ contains an ordered $t$-\toughpair.
In Section~\ref{sec:biclique:nolongpath}, we consider the case when the length of any longest path in $\Dcontr$ is at most $2t+2$. In this case, we get one of the desired outputs.

\subsubsection{When $\Dcontr$ has a long path}\label{sec:biclique:longpath}

In this section, we prove Lemma~\ref{lem:digging} when the length of a longest path in $\Dcontr$ is at least $2t+2$.
Let $P=(v_1, \ldots,v_p)$
be a directed longest path in $\Dcontr$ (then $p \geq 2t+2$).
Recall that none of the edges of $\Dcontr$, and in particular, none  of the edges of $P$, are {\em contraction-redundant} with respect to $(A,B)$.

\begin{lemma}\label{claim:biclique1}
For each $e \in E(P)$, $\Dcontr/e$ is acyclic.
\end{lemma}

\begin{proof}
Suppose, for the sake of contradiction, that there is a cycle in $\Dcontr/e$. Let $e=(v_i,v_{i+1})$. 
Then, either there is a $(v_i,v_{i+1})$-path in $\Dcontr - \{(v_i,v_{i+1})\}$ or a $(v_{i+1},v_i)$-path in $\Dcontr$. 
If there exists a $(v_{i+1},v_i)$-path in $\Dcontr$ then together with the edge $e=(v_i,v_{i+1})$, it creates a directed closed walk in $\Dcontr$. 
This contradicts that $\Dcontr$ is acyclic. 

In the other case, suppose there exists a $(v_i,v_{i+1})$-path, say $P'$ in $\Dcontr -\{(v_i,v_{i+1})\}$. We first claim that the internal vertices of $P'$ are disjoint from $P$. Suppose not, let $v_j \in V(P)$ be the first internal vertex of $P'$ that belongs to $P$. Then there is a $(v_i,v_j)$-path in $\Dcontr$ and $(v_j,v_{i+1})$-path in $\Dcontr$. If $j<i$, then the $(v_i,v_j)$-path together with the $(v_j,v_i)$-subpath of $P$, gives a closed walk in $\Dcontr$, contradicting its acyclicity. Otherwise, $j>i+1$. In this case, the $(v_j,v_{i+1})$-path together with the $(v_{i+1},v_j)$-subpath of $P$ gives a closed walk in $\Dcontr$, again contradicting its acyclicity.
\end{proof}

Since each edge of $\Dcontr$, in particular, each edge of $P$, is not contraction-redundant, from Lemma~\ref{claim:biclique1} and the definition of contraction-redundant, for each $e \in E(P)$ either
\begin{enumerate}
\item there exists an $aa$-path, or a $bb$-path in $\Dcontr/e$ (that is, $(\Dcontr/e)^{\star}$ does not contain the $(A,B)$ -induced-biclique), or 
\item there exists $i \in [g(t)]$ an $a_{i}b_{i}$-path in $(\Dcontr/e )- \{(a_{i},b_{i})\}$ ($(a_{i},b_{i})$ is not a minimal edge in $\Dcontr/e$).
\end{enumerate}

\begin{lemma}\label{obs:biclique1}
Let $e=(v_i,v_{i+1}) \in E(P)$ such that $H_{contr}/e$ has an $(a_j,b_j)$-path. Then there exists an $(a_j,v_{i+1})$-path and a $(v_i,b_j)$-path in $\Dcontr$.

If $\Dcontr/e$ has a $(a_j,a_{\ell})$-path, for some $a_j,a_{\ell} \in A$, then there exists an $(a_j, v_{i+1})$-path and a $(v_i,a_{\ell})$-path in $\Dcontr$.

If $\Dcontr/e$ has a $(b_j,b_{\ell})$-path, for some $b_j,b_{\ell} \in B$, then there exists an $(b_j, v_{i+1})$-path and $(v_i,b_{\ell})$-path in $\Dcontr$.
\end{lemma}

\begin{proof}
Let $e=(v_i,v_{i+1}) \in E(P)$ such that $D_{contr}/e$ has an $(a_j,b_j)$-path. Then either there exists an $(a_j,v_i)$-path and a $(v_{i+1},b_j)$-path in $\Dcontr$, or there exists an $(a_j,v_{i+1})$-path and $(v_i,b_j)$-path in $\Dcontr$. In the later case, we are done. The former case implies an $(a_j,b_j)$-path in $\Dcontr$ of length at least three (with at least three edges), which contradicts that $(a_j,b_j)$ is a minimal edge in $\Dcontr$.

If $D_{contr}/e$ has an $(a_j,a_{\ell})$-path, then either there exists an $(a_j,v_i)$-path and a $(v_{i+1},a_{\ell})$-path in $\Dcontr$, or there exists an $(a_j,v_{i+1})$-path and $(v_i,a_{\ell})$-path in $\Dcontr$. In the later case, we are done. The former case implies an $(a_j,a_{\ell})$-path in $\Dcontr$ which contradicts that $(A,B)$ is an induced biclique in $\Dcontr^{\star}$ (in particular, that $A$ is an independent set in $\Dcontr^{\star}$).
One can similarly show that if $D_{contr}/e$ has a $(b_j,b_{\ell})$-path, then there exists an $(b_j,v_{i+1})$-path and $(v_i,b_{\ell})$-path in $\Dcontr$.
\end{proof}

For each $e=(v_i,v_{i+1}) \in E(P)$, we say that:
\begin{enumerate}
\item $e$ is {\em $j$-irredundant} in $\Dcontr$, if there is an $(a_j,v_{i+1})$-path and a $(v_i,b_j)$-path in $\Dcontr$.
\item $e$ is {\em $aa$-irredundant} in $\Dcontr$, if there exists $a_j,a_{\ell} \in A$ such that there is an $(a_j,v_{i+1})$-path and a $(v_i,a_{\ell})$-path in $\Dcontr$.
\item  $e$ is {\em $bb$-irredundant} in $\Dcontr$, if there exists $b_j,b_{\ell} \in B$ such that there is an $(b_j,v_{i+1})$-path and $(v_i,b_{\ell})$-path in $\Dcontr$.
\end{enumerate}

From Lemmas~\ref{claim:biclique1} and~\ref{obs:biclique1}, for each $e \in E(P)$, either $e$ is $j$-irredundant, for some $j \in [g(t)]$, or is $aa$-irredundant or is $bb$-irredundant.

\begin{lemma}\label{claim:biclique21}
Let $e=(v_i,v_{i+1}) \in E(P)$ such that $e$ is $aa$-irredundant in $\Dcontr$. Then $i=1$.
\end{lemma}

\begin{proof}
For the sake of contradiction, suppose that $i>1$. Let $e'=(v_{i-1},v_i)$. 
If $e'$ is $aa$-irredundant or $e'$ is $j$-irredundant for some $j \in [g(t)]$, then there exists a path from some vertex of $A$ to $v_i$ in $\Dcontr$. Also since $e$ is $aa$-irredundant, there exists a path from $v_i$ to some vertex of $A$ in $\Dcontr$. This implies either a directed closed walk in $\Dcontr$, contradicting its acyclicity, or a path between two distinct vertices of $A$ in $\Dcontr$, contradicting that $A$ is an independent set in $\Dcontr^{\star}$.

If $e'$ is $bb$-irredundant, then there exists a path from some vertex of $B$ to $v_i$. Also since $e$ is $aa$-irredundant, there exists a path from $v_i$ to some vertex of $A$. This implies a path from some vertex of $B$ to some vertex of $A$. Since $(A,B)$ is an induced biclique in $\Dcontr^{\star}$, that is there is a path from every vertex of $A$ to every vertex of $B$ in $\Dcontr$, this implies a directed closed walk in $\Dcontr$, contradicting its acyclicity.
\end{proof}

\begin{lemma}\label{claim:biclique22}
Let $e=(v_i,v_{i+1}) \in E(P)$ such that $e$ is $bb$-irredundant in $\Dcontr$. Then $i+1=p$.
\end{lemma}

\begin{proof}
For the sake of contradiction, suppose that $i+1<p$. Let $e'=(v_{i+1},v_{i+2})$. 
If $e'$ is $bb$-irredundant or $e'$ is $j$-irredundant for some $j \in [g(t)]$, then there exists a path from $v_{i+1}$ to some vertex of $B$. Also since $e$ is $bb$-irredundant, there exists a path from some vertex of $B$ to $v_{i+1}$. This either implies a directed closed walk in $\Dcontr$, contradicting its acyclicity, or a path between two distinct vertices of $B$ in $\Dcontr$, contradicting that $B$ is an independent set in $\Dcontr^{\star}$.

If $e'$ is $aa$-irredundant, then there exists a path from $v_{i+1}$ to some vertex of $A$ in $\Dcontr$. Also since $e$ is $bb$-irredundant, there exists a path from some vertex of $B$ to $v_{i+1}$ in $\Dcontr$. This implies a path from some vertex of $B$ to some vertex of $A$ in $\Dcontr$. Since 
$(A,B)$ is an induced-biclique in $\Dcontr^{\star}$, this implies a directed closed walk
in $\Dcontr$, contradicting its acyclicity.
\end{proof}

Let $P'$ be a subpath obtained from $P$ by without the first and last edge of $P$, that is $P'=(v_2, \ldots,v_{p-1})$. From Lemmas~\ref{claim:biclique21} and~\ref{claim:biclique22}, each edge $e \in E(P')$ is $j$-irredundant for some $j \in [g(t)]$.

\begin{lemma}\label{claim:biclique3}
Let $e,e'$ be distinct edges of $P'$ such that $e$ is $q$-irredundant in $\Dcontr$ and $e'$ is $r$-irredundant in $\Dcontr$, for some $q,r \in [g(t)]$. Then $q \neq r$.
\end{lemma}

\begin{proof}
For the sake of contradiction, suppose that $q=r$. Without loss of generality, let $e=(v_i,v_{i+1})$ and $e'=(v_j,v_{j+1})$ such that $i<j$. Then there exists an $(a_q,v_{i+1})$-path and $(v_j,b_q)$-path in $\Dcontr$. Since $i < j$, there exists a path from $(v_i,v_j)$ (which is the the $(v_{i+1},v_j)$-subpath of $P'$) in $\Dcontr$. This implies an $(a_q,b_q)$-path of length strictly greater than one in $\Dcontr$, contradicting that $(a_q,b_q)$ is a minimal edge in $\Dcontr$.
\end{proof}

For each $e=(v_i,v_{i+1}) \in E(P')$, fix a $\rho(i) \in [g(t)]$, such $e$ is $\rho(i)$-irredundant in $\Dcontr$. 
The following observation follows from Lemma~\ref{claim:biclique3}.

\begin{observation}\label{obs:biclique1}
For each $v_i \in V(P')$, $i \in [3, p-2]$, there exists a $(a_{\rho(i-1)},v_i)$-path and a $(v_i, b_{\rho(i)})$-path in $\Dcontr$.  Also, for $i,j \in [3,p-2]$, $i \neq j$, $\rho(i) \neq \rho(j)$.
\end{observation}

For each $i \in [3,p-2]$, let $P^{\aaa}_i$ denote an (arbitrarily) fixed $(a_{\rho(i-1)},v_i)$-path in $\Dcontr$ and let $P^{\bbb}_i$ denote an (arbitrarily) fixed $(v_i,b_{\rho(i)})$-path in $\Dcontr$.
Let $E_1 = (v_3, b_{\rho(3)}), \allowbreak  (v_4, b_{\rho(4)}), \ldots, \allowbreak (v_{p/2-2}, b_{\rho(p/2-2)})$ and let 
 $E_2 = ((a_{\rho(p/2-2)},v_{p/2-1}), \allowbreak  (a_{\rho(p/2-1)}, \allowbreak v_{p/2}), \ldots, (a_{\rho(p-1)}, \allowbreak v_{p-2})$. 
 From Observation~\ref{obs:biclique1}, each edge of $E_1 \cup E_2$ is an edge in $D^{\star}_{\contr}$ and the endpoints of the edges in $E_1 \cup E_2$ are distinct. We will now show that $(E_1,E_2)$ is an ordered $(p/2-1)$-\toughpair\ in $D^{\star}_{\contr}$. Since $p \geq 2t+2$, this proves Lemma~\ref{lem:digging} in the case when the length of a longest path in $\Dcontr$ is at least $2t+2$.
 
\begin{lemma}\label{lem:long-path-ordered-tough-pair}
$(E_1,E_2)$ is an ordered $(p/2-1)$-\toughpair\ in $D^{\star}_{\contr}$.
\end{lemma}

\begin{proof}
To prove the lemma, we prove the following claims.

\begin{claim}[Weak independence of $E_1$ (resp.~$E_2$)]\label{claim:biclique-long-path-weak}
The edges in $E_1$ are pairwise weakly independent. Similarly, the edges in $E_2$ are pairwise weakly independent in $D^{\star}_{\contr}$.
\end{claim}

\begin{claimproof}
Let $(a_{\rho(i-1)},{v}_{i}), \allowbreak (a_{\rho(j-1)},{v}_{j}) \in E_1$. For the sake of contradiction, say $(v_j, a_{\rho(i-1)}) \in \Dcontr^{\star}$. Then this implies a $(a_{\rho(j-1},a_{\rho(i-1)})$-path in $\Dcontr^{\star}$ contradicting that $A$ is an independent set in $\Dcontr^{\star}$.
Using symmetric arguments one can show that all edges in $E_2$ are pairwise weakly independent.
\end{claimproof}

\begin{claim}[Orderedness on $E_1$ and $E_2$]\label{claim:biclique-long-path-ordered}
For any $i,j \in [3, p-2]$, $i <j$, $(a_{\rho(j-1)}, v_{i}), \allowbreak (a_{\rho(j-1)}, a_{\rho(i-1)}),  \allowbreak ({v}_{j},  \allowbreak b_{\rho(i)}), \allowbreak (v_j,v_i) \not \in E(\Dcontr^{\star})$.
\end{claim}

\begin{claimproof}
Since $A$ is an independent set in $\Dcontr^{\star}$, $(a_{\rho(j-1)}, a_{\rho(i-1)}) \not \in E(\Dcontr^{\star})$. Since $\Dcontr^{\star}$ is acyclic and there is a $(v_i,v_j)$-path in $\Dcontr$, we conclude that $(v_j,v_i) \not \in E(\Dcontr^{\star})$.
If there is an $(a_{\rho(j-1)}, {v}_{i})$-path in $\Dcontr$, then this together with the $( {v}_{i},v_{j-1})$-subpath of $P'$ and the $(v_{i-1},b_{\rho(i-1)})$ path $P^{\bbb}_{i-1}$, implies a $(a_{\rho(j-1),b_{\rho(j-1)}})$-path of length strictly greater than $1$ in $\Dcontr$, which is a contradiction. Using symmetric arguments, one can show that 
$({v}_{j}, b_{\rho(i)}), (v_j,v_i) \not \in E(\Dcontr^{\star})$.
\end{claimproof}

\begin{claim}[Weak independence between $E_1$ and $E_2$]\label{claim:biclique-long-path-weak-between}
Let $(v_i,b_{\rho(i)}) \in E_1$ and $(a_{\rho(j-1)},v_j) \in E_2$.
There is $(v_j,b_{\rho(i)}), (b_{\rho(i)},a_{\rho(i-1)}) \not \in E(\Dcontr^{\star})$. 
\end{claim}

\begin{claimproof}
First observe from the construction of $E_1,E_2$ that $i <j$.
If $(v_j,b_{\rho(i)}) \in E(\Dcontr^{\star})$, then the $(a_{\rho(i)},v_{i+1})$-path $P^{\aaa}_{i+1}$, together with the $(v_{i+1},v_j)$-subpath of $P'$, together with the $(v_j,b_{\rho(i)})$-path, implies a $(a_{\rho(i)},b_{\rho(i)})$-path of length strictly greater than $1$ in $\Dcontr$, which is a contradiction. 

If $(b_{\rho(i)},a_{\rho(i-1)})  \in E(\Dcontr^{\star})$, then the $(a_{rho(i-1)},b_{\rho(i)})$-path obtained by concatenating $P^{\aaa}_i$ and $P^{\bbb}_i$, together with the $(b_{\rho(i)},a_{\rho(i-1)})$-path, implies a $(a_{\rho(i-1)},a_{\rho(j-1)})$-path in $\Dcontr$, which contradicts that $A$ is an independent set in $\Dcontr$.
\end{claimproof}

\begin{claim}[Strong independence between $E_1$ and $E_2$]\label{claim:biclique-long-path-strong}
Let $(v_i,b_{\rho(i)}) \in E_1$ and $(a_{\rho(j-1)},v_j) \in E_2$.
Then $(v_i,a_{\rho(j-1}), \allowbreak (a_{\rho(j-1)},v_i), \allowbreak (b_{\rho(i)},v_j), \allowbreak (v_j,b_{\rho(i)}) \not \in E(\Dcontr^{\star})$.
\end{claim}

\begin{claimproof}
Recall that $i >j$.
If $(v_i,a_{\rho(j-1)}) \in E(\Dcontr^{\star})$, then together with $P^{\aaa}_i$, it contradicts that $A$ is independent in $\Dcontr^{\star}$. Similarly, if $(b_{\rho(i)},v_j) \in E(\Dcontr^{\star})$, then it contradicts the independence of $B$ in $\Dcontr^{\star}$.

If $(a_{\rho(j-1)},v_i) \in E(\Dcontr^{\star})$, then this together with the $(v_i,v_{j-1})$-subpath of $P'$, together with the path $P^{\bbb}_{j-1}$, implies a $(a_{\rho(j-1)},b_{\rho(j-1)})$-path of length strictly greater than $1$ in $\Dcontr$, which is a contradiction. Similarly if, $(v_j,b_{\rho(i)}) \in E(\Dcontr^{\star})$, then this implies a $(a_{\rho(i)},b_{\rho(i)})$-path of length strictly greater than $1$ in $\Dcontr$, which is a contradiction.
\end{claimproof}
This concludes the proof.
\end{proof}

\subsubsection{No long path in $\Dcontr$}\label{sec:biclique:nolongpath}

In this section, we prove Lemma~\ref{lem:digging} in the case when the length of a longest path in $\Dcontr$ is $p \leq 2t+2$. Recall that $\Dcontr$ contains $(A,B)$-matching and $\Dcontr^{\star}$ contains $(A,B)$-induced-biclique. Also, the edges of $A \cup B$ are minimal in $\Dcontr$. We prove Lemma~\ref{lem:digging} by induction on the length of a longest $(A,B)$-path in $\Dcontr$.
For this purpose we essentially restate Lemma~\ref{lem:digging} in a form that is ``induction-friendly''. Observe, as a base case, that when the length of a longest $(A,B)$-path is $1$, then since $\Dcontr^{\star}$ contains $(A,B)$-induced-biclique, we conclude that $\Dcontr$ contains $(A,B)$-induced-biclique and the edges of the $(A,B)$-induced-biclique are minimal. In this case, we conclude that we get the second output of Lemma~\ref{lem:digging}.

\begin{lemma}\label{lem:biclique:nolongpath}
Let $A,B \subseteq V(\Dcontr)$ be ordered sets such that the length
of a longest $(A,B)$-path in $\Dcontr$ is $p$,
$|A|= |B| \geq f(t,p)$, $\Dcontr^{\star}$ contains an $(A,B)$-induced-biclique and the matching edges of $(A,B)$ are minimal in $\Dcontr$, then either
\begin{enumerate}
\item there exist ordered sets $A',B' \subseteq V(\Dcontr)$ such that the length
of a longest $(A',B')$-path in $\Dcontr$ is at most $p$,
$|A'|= |B'| \geq f(t,p-1)$, $\Dcontr^{\star}$ contains an $(A',B')$-induced-biclique and the matching edges of $(A',B')$ are minimal in $\Dcontr$, or
\item $\Dcontr$ contains an ordered $t$-\toughpair, or
\item $\Dcontr^{\star}$ contains a $t$-induced-biclique whose edges are minimal in $\Dcontr$.
\end{enumerate}
\end{lemma}

\begin{proof}
Let $L$ be the set of vertices of $\Dcontr$ that contains the first internal vertex on every longest $(A,B)$-path in $\Dcontr$.

\begin{claim}\label{claim:biclique6}
There is no $(L,L)$-path in $\Dcontr$.
\end{claim}

\begin{claimproof}
For the sake of contradiction, suppose there exist $x,y \in L$ such that there this an $(x,y)$-path in $\Dcontr$. Note that $x\neq y$, as otherwise, $\Dcontr$ would contain a cycle.
From the definition of $L$, there exists $a_i,a_j \in L$ ($i$ not necessarily distinct from $j$) such that $x$ is the first internal vertex on some $(a_i,B)$-path, say $P_x$, of length $p$ and $y$ is the first internal vertex on some $(a_j,B)$-path, say $P_y$, of length $p$. 
Let $P_{xy}$ be some $(x,y)$-path in $\Dcontr$. We first claim that the set of internal vertices of $P_{x,y}$ is disjoint from $P_y$. Suppose not. Let $z$ be some internal vertex of $P_{xy}$ that is also a vertex of $P_y$. Since $z \in V(P_y)$, there exists a $(y,z)$-path in $\Dcontr$. Also since $z \in V(P_{x,y})$, there exists a $(z,y)$-path in $\Dcontr$. This implies a cycle in $\Dcontr$, which is a contradiction. Thus, we conclude that the internal vertices of $P_{xy}$ are disjoint from that of $P_y$. Then consider the $(x,B)$-path obtained by appending $P_{xy}$ and $P_y$. Note that the length of $P_y$ is $p-1$ and the length of $P_{xy}$ is at least $1$. This $(x,B)$-path, together with the edge $(a_i,x)$, gives an $(a_i,B)$-path of length strictly greater than $p$ in $\Dcontr$, which contradicts that the length of a longest $(A,B)$-path in $\Dcontr$ is $p$.
\end{claimproof}

\begin{claim}\label{claim:biclique7}
For each $(a,x) \in E(\Dcontr)$ such that $a \in A$ and $x \in L$, $(a,x)$ is a minimal edge of $\Dcontr$. 
\end{claim}

\begin{claimproof}
For the sake of contradiction, say there exists an $(a,x)$-path in $\Dcontr$ of length at least $2$. Let this path be $P_{ax}$. Since $x \in L$, there exists an $(A,B)$-path of length $p$ whose first internal vertex is $x$. Let this path be $P_{AB}$. Then the internal vertices of $P_{ax}$ are disjoint from the internal vertices of $P_{AB}$, as otherwise there would be a cycle in $\Dcontr$. By appending the path $P_{sx}$ (which is of length at least $2$) with the subpath of $P_{ABs}$ starting from $x$ (which is of length $p-1$), we get an $(A,B)$-path in $\Dcontr$ of length at least $p+1$, which is a contradiction.
\end{claimproof}

\begin{claim}\label{claim:biclique8}
Let $a \in A$ and $x \in L$, such that $(a,x) \not \in E(\Dcontr)$. Then $(a,x) \not \in E(\Dcontr^{\star})$.
\end{claim}

\begin{claimproof}
For the sake of contradiction, suppose there exists an $(a,x)$-path, say $P_{a,x}$ in $\Dcontr$ of length at least two. Since $x \in L$, there exists $a' \in A$ ($a'$ not necessarily distinct from $a$) such that $x$ is the first internal vertex of a $(a',T)$-path, say $P_{a'B}$, of length $p$. First observe that the internal vertices of $P_{a,x}$ are disjoint from that of $P_{a'B}$, as otherwise there would be a cycle in $\Dcontr$. The path $P_{a,x}$ appended with the $(x,B)$-subpath of $P_{a'B}$ gives an $(A,B)$-path in $\Dcontr$ of length strictly greater than $p$, which is a contradiction.
\end{claimproof}

Consider the bipartite graph $D_{\bip}=\Dcontr[A \cup L]$. Note, from Claim~\ref{claim:biclique7}, that each edge of $D_{\bip}$ is a minimal edge in $\Dcontr$.
We now distinguish into two cases based on the size of the matching in $D_{\bip}$.\\

\noindent{\bf Case 1: The size of a maximum matching in $D_{\bip}$ is at least $\mathcal{R}'(2t)$.} In this case, from Ramsey's Theorem, either there exists an induced matching of size $2t$ in $D_{\bip}$ or a $2t$-induced-biclique. From Claim~\ref{claim:biclique7}, each edge of $D_{\bip}$ is a minimal edge of $D_{contr}$. 

In the first case, we get sets $A^* \subseteq A$ and $B^* \subseteq L$ such that there is an $(A^*,B^*)$-induced-matching in $D_{\bip}$. Since $A$ is an independent set in $\Dcontr^{\star}$, so is $A^*$. Since $B^* \subseteq L$, from Claim~\ref{claim:biclique7}, $B^*$ is an independent set in $\Dcontr^{\star}$. Thus, $(A^*,B^*)$ is an induced-matching in $\Dcontr$. Further, from Claim~\ref{claim:biclique8}, $(A^*,B^*)$ is in fact an induced matching in $\Dcontr^{\star}$. Observe that if a digraph has an induced-matching of size $2t$ in its transitive closure, then it has a ordered $t$-\toughpair. Thus, in this case we get the second outcome of the lemma.

In the second case, when $H_{\bip}$ contains a $2t$-induced-biclique, then this is also a $2t$-induced-biclique in $D_{\contr}$ where all edges of the biclique are minimal edges of $\Dcontr$. Thus, in this case we get the third outcome of the lemma.\\

\noindent{\bf Case 2: The size of a maximum matching in $D_{\bip}$ is at most $\mathcal{R}'(2t)$.} In this case, 
we will find ordered sets $A',B' \subseteq V(\Dcontr)$ satisfies the properties stated in the first outcome of the lemma.

Let $Z$ be a minimum vertex cover of $D_{\bip}$ of size at most $\mathcal{R}'(2t)$. Let $Z_A = Z \cap A$ and 
$Z_L=Z \cap L$. Let $A^*=A \setminus Z_S$. Since $Z_A \cup Z_L$ is a vertex cover of $D_{\bip}$, for any  $a \in A^*$, $N_L(a) \subseteq Z_L$.
For each subset $L' \subseteq Z_L$, let $A_{L'} \subseteq A^*$ such that for each vertex $a \in A_{L'}$, $N_{Z_L}(a)=L'$ (and hence $N_L(a)=L'$).

Fix $L' \subseteq Z_L$ such that $|S_{L'}|$ is maximized. 
Recall that the sets $A =(a_1,\ldots,a_{f(t,p)})$ and $T=(t_1, \ldots, t_{f(t,p)})$ are ordered sets.
Set $A'=A_{L'}$ and $B'$ be the corresponding vertices of $B$, 
that is $b_i \in B'$ if and only if $a_i \in A'$.
From the choice of $L'$, $|A'|=|B'| \geq (|A|-\mathcal{R}'(2t))/2^{\mathcal{R}'(2t)} = f(t,p-1)$.
We will now show that the length of a longest $(A',B')$-path in $\Dcontr$ is at most $p-1$. 

\begin{claim}\label{claim:biclique5}
Let $a_i \in A'$ and $b_j \in B'$ such that $i \neq j$. Let $P=(a_i,x_1,\ldots,x_q,b_j)$ be a longest $(A',B')$-path in $\Dcontr$. Then $x_1 \not \in L$.
\end{claim}
\begin{claimproof}
Since $b_j \in B'$, $a_j \in A'$ (from construction of $B'$).
First observe that since $b_j \in B'$, then $a_j \in A'$.
For the sake of contradiction, suppose $x_1 \in L$. Since $(a_i,x) \in E(\Dcontr)$ and all the vertices of $A'$ have the same neighbourhood in $Z_L$ and $a_j \in A'$, we conclude that $(a_j,x) \in E(\Dcontr)$. Thus, there is an $(a_j,b_j)$-path in $\Dcontr$ of length strictly greater than one, which contradicts that $(a_j,b_j)$ is a minimal edge of $\Dcontr$.
\end{claimproof}

From Claim~\ref{claim:biclique5}, the first internal vertex of every longest $(A',B')$-path is not contained in $L$. Now suppose that there exists an $(A',B')$-path of length $p$. Then its first internal vertex should be in $L$ by the definition of $L$. This is a contradiction.  
\end{proof}

Lemma~\ref{lem:biclique:nolongpath} proves Lemma~\ref{lem:digging} when the length of the longest path $p \leq 2t+2$.